\documentclass[a4paper,11pt]{article}

% ============================================================
%  Le Collé — Séance 1 : Déterminants, courbes paramétrées
%  © 2026 Ahmed Amine Ouqlani
% ============================================================

% ---------- Polices ----------
\usepackage[utf8]{inputenc}
\usepackage[T1]{fontenc}
\usepackage[sc]{mathpazo}          % texte et mathématiques : Palatino
\usepackage[scaled=0.88]{helvet}   % titres et repères : linéale
\linespread{1.07}
\usepackage{amsmath,amssymb,amsthm,mathrsfs,stmaryrd}
\usepackage{microtype}

% ---------- Mise en page ----------
\usepackage[a4paper,top=2.3cm,bottom=2.3cm,left=2.4cm,right=2.4cm,
            headheight=14pt,headsep=0.8cm,footskip=1.1cm]{geometry}
\setlength{\parindent}{0pt}
\setlength{\parskip}{0.45em plus 0.15em minus 0.1em}
\raggedbottom
\emergencystretch=1.5em
\clubpenalty=10000 \widowpenalty=10000 \displaywidowpenalty=10000 \brokenpenalty=10000

% ---------- Couleurs : un bleu de nuit, un ocre, des gris ----------
\usepackage{xcolor}
\definecolor{Ink}{HTML}{1C2533}
\definecolor{Bleu}{HTML}{1F3A5F}
\definecolor{BleuMoyen}{HTML}{4A6A8F}
\definecolor{BleuClair}{HTML}{9DB2C9}
\definecolor{Fond}{HTML}{F2F4F7}
\definecolor{Ocre}{HTML}{B07A26}
\definecolor{Brique}{HTML}{A23B2A}
\definecolor{Gris}{HTML}{5E6672}
\definecolor{Filet}{HTML}{C5CCD6}

% ---------- Packages ----------
\usepackage{enumitem}
\setlist{itemsep=0.25em,topsep=0.35em,parsep=0pt}
\setlist[itemize]{label=\textcolor{BleuMoyen}{\small\textbullet},leftmargin=1.4em}
\setlist[enumerate]{leftmargin=1.8em,label=\textsf{\small\color{Bleu}\arabic*.}}
\usepackage{booktabs,array,multicol}
\usepackage{needspace,float}
\usepackage[font=small,labelfont={sf,color=Bleu},labelsep=period,skip=6pt]{caption}
\usepackage{tikz}
\usetikzlibrary{arrows.meta,calc,decorations.markings}
\usepackage{pgfplots}
\pgfplotsset{compat=1.17}
\usepackage{tkz-tab}
\usepackage[most]{tcolorbox}
\usepackage{titlesec,titletoc}
\usepackage{fancyhdr,lastpage,etoolbox}
\usepackage[colorlinks=true,linkcolor=Bleu,urlcolor=Bleu,citecolor=Bleu,
            pdftitle={Le Collé — Séries numériques : sélection de 5 problèmes},pdfauthor={Ahmed Amine Ouqlani}]{hyperref}

\numberwithin{equation}{section}
\numberwithin{figure}{section}
\numberwithin{table}{section}

% ---------- Titres ----------
% Chapitres (sections) : grand numéro, titre en linéale, filet ocre
\titleformat{\section}[display]
  {\normalfont\sffamily\color{Ink}}
  {{\small\color{Gris}\textls[200]{\MakeUppercase{Chapitre}}}\\[-0.2em]%
   {\fontsize{64}{64}\selectfont\color{BleuClair}\thesection}}
  {0.4em}
  {\fontsize{25}{30}\selectfont\bfseries}
  [\vspace{0.5em}{\color{Ocre}\rule{3.2cm}{1.4pt}}]
\titlespacing*{\section}{0pt}{0pt}{1.6em}
\titleformat{name=\section,numberless}[display]
  {\normalfont\sffamily\color{Ink}}{}{0pt}{\fontsize{25}{30}\selectfont\bfseries}
  [\vspace{0.5em}{\color{Ocre}\rule{3.2cm}{1.4pt}}]
\titlespacing*{name=\section,numberless}{0pt}{0pt}{1.6em}

\titleformat{\subsection}
  {\needspace{5\baselineskip}\normalfont\sffamily\Large\bfseries\color{Ink}}
  {\color{Bleu}\thesubsection}{0.8em}{}
\titlespacing*{\subsection}{0pt}{1.8em}{0.5em}
\titleformat{\subsubsection}
  {\needspace{5\baselineskip}\normalfont\sffamily\large\bfseries\color{Ink}}
  {\color{Bleu}\thesubsubsection}{0.7em}{}
\titlespacing*{\subsubsection}{0pt}{1.6em}{0.4em}
\titleformat{\paragraph}[runin]{\normalfont\sffamily\bfseries\color{Ink}}{}{0pt}{}[.]
\titlespacing*{\paragraph}{0pt}{0.9em}{0.6em}
\setcounter{secnumdepth}{3}
\setcounter{tocdepth}{2}

% ---------- Table des matières ----------
\titlecontents{section}[0pt]{\addvspace{0.7em}\sffamily\bfseries\color{Ink}}
  {\makebox[2.2em][l]{\color{Bleu}\thecontentslabel}}{}{\hfill\contentspage}[\addvspace{0.2em}]
\titlecontents{subsection}[2.2em]{\small\linespread{1}\selectfont}
  {\makebox[2.6em][l]{\color{Gris}\thecontentslabel}}{}{\titlerule*[0.7pc]{.}\contentspage}

% ---------- En-têtes et pieds de page ----------
\pagestyle{fancy}
\fancyhf{}
\renewcommand{\sectionmark}[1]{\markboth{\thesection\quad#1}{}}
\fancyhead[L]{\footnotesize\sffamily\color{Gris}\leftmark}
\fancyhead[R]{\footnotesize\sffamily\scshape\color{Gris}Le Collé}
\fancyfoot[C]{\small\sffamily\color{Ink}\thepage}
\fancyfoot[L]{\scriptsize\sffamily\hyperlink{toc}{\color{Filet}sommaire}}
\renewcommand{\headrulewidth}{0.4pt}
\renewcommand{\headrule}{{\color{Filet}\hrule height\headrulewidth width\headwidth}}
\renewcommand{\footrulewidth}{0pt}
\fancypagestyle{plain}{\fancyhf{}\fancyfoot[C]{\small\sffamily\color{Ink}\thepage}\renewcommand{\headrulewidth}{0pt}}

% ---------- Énoncés ----------
\newtheoremstyle{enonce}{0pt}{0pt}{}{}{\sffamily\bfseries\color{Bleu}}{.}{0.5em}
  {\thmname{#1}\thmnumber{~#2}\thmnote{\ \normalfont\sffamily\color{Ink}(#3)}}
\newtheoremstyle{exemple}{0.9em}{0.9em}{}{}{\sffamily\bfseries\color{BleuMoyen}}{.}{0.5em}
  {\thmname{#1}\thmnumber{~#2}\thmnote{\ \normalfont\sffamily\color{Ink}(#3)}}
\newtheoremstyle{remarque}{0.7em}{0.7em}{}{}{\sffamily\itshape}{.}{0.5em}{\thmname{#1}\thmnote{\ (#3)}}
\theoremstyle{enonce}
\newtheorem{definition}{Définition}[section]
\newtheorem{theoreme}[definition]{Théorème}
\newtheorem{proposition}[definition]{Proposition}
\newtheorem{corollaire}[definition]{Corollaire}
\theoremstyle{exemple}
\newtheorem{exemple}[definition]{Exemple}
\theoremstyle{remarque}
\newtheorem*{remarque}{Remarque}

\tcbset{
  defi/.style={enhanced,breakable,colback=white,colframe=white,boxrule=0pt,arc=0pt,
     borderline west={1.2pt}{0pt}{Bleu},left=11pt,right=4pt,top=3pt,bottom=3pt,
     before skip=1.1em,after skip=1.1em},
  thm/.style={enhanced,breakable,colback=Fond,colframe=Fond,boxrule=0pt,arc=0pt,
     borderline west={2.4pt}{0pt}{Bleu},left=11pt,right=8pt,top=6pt,bottom=6pt,
     before skip=1.2em,after skip=1.2em}}
\tcolorboxenvironment{definition}{defi}
\tcolorboxenvironment{theoreme}{thm}
\tcolorboxenvironment{proposition}{thm}
\tcolorboxenvironment{corollaire}{thm}

\renewcommand{\proofname}{Démonstration}
\renewcommand{\qedsymbol}{\textcolor{BleuClair}{$\blacksquare$}}
\makeatletter
\renewenvironment{proof}[1][\proofname]{\par\pushQED{\qed}\normalfont
  \topsep0.5em\relax\trivlist\item[\hskip\labelsep\sffamily\itshape\color{Gris}#1.]\ignorespaces}
  {\popQED\endtrivlist\@endpefalse}
\makeatother

% ---------- Méthode : filets haut et bas, titre en bandeau ----------
\makeatletter\newcommand{\methlabelfix}{\protected@edef\@currentlabel{\thetcbcounter}}\makeatother
\newtcolorbox[auto counter,number within=section]{methode}[2][]{%
  enhanced,unbreakable,colback=white,colframe=white,boxrule=0pt,arc=0pt,
  borderline north={0.9pt}{0pt}{Bleu},borderline south={0.9pt}{0pt}{Bleu},
  left=4pt,right=4pt,top=8pt,bottom=8pt,before skip=1.4em,after skip=1.4em,
  before upper={\methlabelfix\phantomsection\addcontentsline{met}{metentry}{\protect\makebox[2.6em][l]{\thetcbcounter}#2}%
     {\sffamily\bfseries\small\color{white}\colorbox{Bleu}{\strut\ MÉTHODE~\thetcbcounter\ }}%
     \ \ {\sffamily\bfseries\color{Bleu}#2}\par\smallskip},#1}

% ---------- Repères pédagogiques : un filet de couleur, un mot-clé ----------
\newcommand{\definirbloc}[3]{% nom, couleur, libellé
  \newtcolorbox{#1}[1][]{enhanced,breakable,blanker,
    borderline west={1.4pt}{0pt}{#2},left=10pt,right=2pt,top=1pt,bottom=1pt,
    before skip=1em,after skip=1em,
    before upper={{\sffamily\bfseries\footnotesize\color{#2}\textls[120]{\MakeUppercase{#3}}}\quad},##1}}
\definirbloc{reflexe}{Bleu}{Réflexe}
\definirbloc{astuce}{Ocre}{Astuce}
\definirbloc{attention}{Brique}{Attention}
\definirbloc{enpratique}{BleuMoyen}{En pratique}
\definirbloc{plusloin}{Gris}{Pour aller plus loin}
\definirbloc{strategie}{Bleu}{Stratégie}
\definirbloc{lecture}{BleuMoyen}{Lire le schéma}
% compatibilité avec les parties reprises
\newenvironment{garde}{\begin{attention}}{\end{attention}}
\newenvironment{complementHP}{\begin{plusloin}}{\end{plusloin}}
\newenvironment{interpretation}{\begin{enpratique}}{\end{enpratique}}

% ---------- Exercices, corrigés ----------
\newcounter{exo}[section]
\renewcommand{\theexo}{\thesection.\arabic{exo}}
\newcommand{\niveau}[1]{{\color{Ocre}\footnotesize
  \ifcase#1\or$\bullet\circ\circ\circ\circ$\or$\bullet\bullet\circ\circ\circ$\or$\bullet\bullet\bullet\circ\circ$%
  \or$\bullet\bullet\bullet\bullet\circ$\or$\bullet\bullet\bullet\bullet\bullet$\fi}}
\newcommand{\serie}[1]{\par\needspace{8\baselineskip}\vspace{1.5em}%
  \noindent{\sffamily\bfseries\color{Bleu}#1}\par\nobreak\vspace{-0.4em}%
  \noindent{\color{Filet}\rule{\linewidth}{0.4pt}}\par\nobreak\vspace{0.1em}}
\newenvironment{exo}[3]{%
  \par\needspace{5\baselineskip}\vspace{0.75em}\refstepcounter{exo}\label{e:#1}%
  \def\exolab{#1}%
  \noindent{\sffamily\bfseries\color{Bleu}Exercice~\theexo}\quad\niveau{#2}%
  \hfill{\footnotesize\sffamily\color{Gris}#3}\par\nobreak\vspace{0.15em}}%
  {\par\nopagebreak\noindent\hfill{\scriptsize\sffamily\hyperref[c:\exolab]{\color{Gris}corrigé p.~\pageref{c:\exolab}}}\par}
\newenvironment{corrige}[1]{%
  \par\needspace{6\baselineskip}\vspace{1em}\phantomsection\label{c:#1}%
  \noindent{\sffamily\bfseries\color{Bleu}Exercice~\ref{e:#1}}\quad
  {\scriptsize\sffamily\hyperref[e:#1]{\color{Gris}énoncé p.~\pageref{e:#1}}}\par\nobreak\vspace{0.2em}}%
  {\hfill\textcolor{BleuClair}{$\square$}\par}
\newcommand{\etape}[1]{\par\smallskip\noindent{\sffamily\bfseries\small\color{BleuMoyen}#1.}\ }
\newenvironment{solution}{\par\noindent{\sffamily\bfseries\small\color{BleuMoyen}Solution.}\ }{\hfill\textcolor{BleuClair}{$\square$}\par}
\newcommand{\chercher}{\par\smallskip\noindent{\color{Filet}\rule[0.5ex]{1.2cm}{0.4pt}}\
  {\footnotesize\sffamily\itshape\color{Gris}laisser chercher l'étudiant avant de poursuivre}\par\smallskip}

% ---------- Guide de séance ----------

\newcommand{\marge}[1]{\leavevmode\marginpar{\raggedleft\scriptsize\sffamily\color{Gris}\textls[100]{\MakeUppercase{#1}}}\ignorespaces}
\newcommand{\Etape}[4]{%
  \subsection[Étape #1 --- #3 {\color{Gris}(#2)}]{#3}\label{etape:#1}%
  \noindent{\small\sffamily\color{Gris}{\color{Ocre}\bfseries Étape #1}\quad·\quad #2\quad·\quad #4}\par\medskip}
\newcommand{\question}[1]{\par\smallskip\noindent{\sffamily\bfseries\small\color{BleuMoyen}Question.}\ \textit{#1}\par\smallskip}
\newcommand{\rappel}[1]{\par\smallskip\noindent\phantomsection\addcontentsline{rap}{rapentry}{#1}%
  {\sffamily\bfseries\scriptsize\colorbox{Fond}{\color{Gris}\strut RAPPEL}}\ {\sffamily\small\color{Gris}#1.}\ }
\newcommand{\pourquoi}{\par\noindent{\sffamily\bfseries\small\color{Ocre}Pourquoi ?}\ }

% Ouverture de chapitre : \ouverture{objectifs}{prérequis}{durée}
\newcommand{\ouverture}[3]{%
  \begin{tcolorbox}[enhanced,colback=Fond,colframe=Fond,arc=0pt,boxrule=0pt,left=12pt,right=12pt,top=10pt,bottom=10pt,
     before skip=1.2em,after skip=1.6em]
  \begin{minipage}[t]{0.56\linewidth}
  {\sffamily\bfseries\small\color{Bleu}OBJECTIFS}\par\vspace{0.3em}\small #1
  \end{minipage}\hfill
  \begin{minipage}[t]{0.38\linewidth}
  {\sffamily\bfseries\small\color{Bleu}PRÉREQUIS}\par\vspace{0.3em}\small #2
  \ifx\relax#3\relax\else\par\vspace{0.6em}{\sffamily\bfseries\small\color{Bleu}DURÉE}\par\vspace{0.2em}\small #3\fi
  \end{minipage}
  \end{tcolorbox}}
% Fin de chapitre
\newenvironment{essentiel}{\par\needspace{12\baselineskip}\vspace{1.5em}%
  \begin{tcolorbox}[enhanced,unbreakable,colback=white,colframe=white,boxrule=0pt,arc=0pt,
   borderline north={1.4pt}{0pt}{Ocre},left=2pt,right=2pt,top=8pt,bottom=4pt]
  {\sffamily\bfseries\color{Ocre}L'ESSENTIEL DU CHAPITRE}\par\medskip}{\end{tcolorbox}}

% ---------- Notations ----------
\newcommand{\K}{\mathbb{K}}
\newcommand{\R}{\mathbb{R}}
\newcommand{\C}{\mathbb{C}}
\newcommand{\N}{\mathbb{N}}
\newcommand{\Z}{\mathbb{Z}}
\newcommand{\Mn}[1]{\mathcal{M}_{n}(#1)}
\newcommand{\GLn}[1]{\mathrm{GL}_n(#1)}
\newcommand{\LE}{\mathcal{L}(E)}
\newcommand{\rg}{\operatorname{rg}}
\newcommand{\vect}{\operatorname{Vect}}
\newcommand{\Mat}{\operatorname{Mat}}
\newcommand{\Id}{\operatorname{Id}}
\newcommand{\sh}{\operatorname{sh}}
\newcommand{\ch}{\operatorname{ch}}
\renewcommand{\th}{\operatorname{th}}
\DeclareMathOperator{\argsh}{argsh}
\newcommand{\vv}[1]{\overrightarrow{#1}}
\newcommand{\ve}[1]{\vec{#1}}
\newcommand{\ient}[2]{\llbracket #1,#2\rrbracket}
\newcommand{\dt}[1]{\left|\begin{matrix}#1\end{matrix}\right|}
\newcommand{\dis}{\displaystyle}

% ---------- Figures ----------
\pgfplotsset{trig format plots=rad}
\pgfplotsset{courbe/.style={axis lines=middle,axis equal image,
  axis line style={Gris,-{Stealth[length=5pt]}},
  xlabel={$x$},ylabel={$y$},ticklabel style={font=\scriptsize\sffamily,color=Gris},
  every axis x label/.style={at={(ticklabel* cs:1)},anchor=west,font=\small},
  every axis y label/.style={at={(ticklabel* cs:1)},anchor=south,font=\small},
  samples=300,smooth,width=0.6\linewidth}}
\tikzset{arcetude/.style={Bleu,line width=1.25pt},
         reste/.style={BleuClair,line width=1pt},
         tgte/.style={Ocre,line width=0.9pt},
         asy/.style={Gris,dashed,line width=0.6pt},
         pt/.style={circle,fill=Ink,inner sep=1.2pt}}

% ---------- Navigation ----------
\makeatletter
\newcommand*\l@rapentry{\@dottedtocline{1}{0em}{0em}}
\newcommand*\l@metentry{\@dottedtocline{1}{0em}{0em}}
\newcommand{\listemethodes}{\@starttoc{met}}
\newcommand{\listerappels}{\@starttoc{rap}}
\makeatother
% onglets de chapitre sur le bord droit
\usepackage{eso-pic}
\newif\ifonglets \ongletsfalse
\newcommand{\dessineronglet}{%
  \edef\ongcol{\ifnum\value{section}=3 Ocre\else Bleu\fi}%
  \edef\ongnum{\the\value{section}}%
  \pgfmathsetmacro{\ongy}{3.2+2.3*(\ongnum-1)}%
  \begin{tikzpicture}[remember picture,overlay]
  \fill[\ongcol] ([xshift=-0.75cm,yshift=-\ongy cm]current page.north east)
     rectangle ([yshift=-\ongy cm-1.9cm]current page.north east);
  \node[text=white,font=\sffamily\bfseries\large] at ([xshift=-0.37cm,yshift=-\ongy cm-0.95cm]current page.north east)
     {\ongnum};
  \end{tikzpicture}}
\AddToShipoutPictureBG{\ifonglets\ifnum\value{section}>0 \dessineronglet\fi\fi}
% blocs de navigation et d'enseignement
\definirbloc{commencer}{Gris}{Si vous commencez ici}
\definirbloc{horsprogramme}{Brique}{Hors programme}
\newcommand{\HP}{\,{\sffamily\bfseries\scriptsize\textcolor{Brique}{[HP]}}}
\newcommand{\source}[1]{{\footnotesize\sffamily\color{Gris}\,(#1)}}
\newcommand{\coupdepouce}[1]{\par\vspace{-0.2em}\noindent{\footnotesize\sffamily\color{Gris}%
  \textbf{Coup de pouce} (à donner si l'étudiant bloque) : #1}\par}
\newcommand{\pointarret}[1]{\par\vspace{1em}\noindent{\color{Filet}\rule[0.5ex]{\linewidth}{0.4pt}}\par\nobreak
  \noindent{\footnotesize\sffamily\color{Ocre}\bfseries POINT D'ARRÊT POSSIBLE}\quad{\footnotesize\sffamily\color{Gris}#1}\par\vspace{0.6em}}
\newcommand{\demarche}[1]{\noindent{\sffamily\bfseries\small\color{BleuMoyen}Démarche.}\ \textit{#1}\par\smallskip}

% ---------- Corrigés : idées, astuces, pièges, à retenir ----------
\newcommand{\choix}[1]{\par\noindent{\sffamily\bfseries\small\color{BleuMoyen}Choix de la méthode.}\ #1\par\smallskip}
\newcommand{\gain}[1]{\par\noindent{\sffamily\bfseries\small\color{Ocre}Astuce.}\ #1\par}
\newcommand{\piege}[1]{\par\noindent{\sffamily\bfseries\small\color{Brique}Piège.}\ #1\par}
\newtcolorbox{retenirbox}{enhanced,breakable,colback=Fond,colframe=Fond,arc=0pt,boxrule=0pt,
  borderline west={2pt}{0pt}{Ocre},left=8pt,right=6pt,top=4pt,bottom=4pt,before skip=0.7em,after skip=0.5em}
\newcommand{\retenir}[1]{\begin{retenirbox}{\sffamily\bfseries\footnotesize\color{Ocre}À RETENIR}\quad #1\end{retenirbox}}
\newcommand{\TD}[1]{{\color{Ocre}\bfseries TD\,#1}\ ·\ }
\newcommand{\q}[1]{\par\smallskip\noindent{\sffamily\bfseries\color{Bleu}#1}\ }
\providecommand{\iddots}{\mathinner{\mkern1mu\raise1pt\hbox{.}\mkern2mu\raise4pt\hbox{.}\mkern2mu\raise7pt\hbox{.}\mkern1mu}}

\providecommand{\Ker}{\operatorname{Ker}}\renewcommand{\Im}{\operatorname{Im}}\providecommand{\tr}{\operatorname{tr}}\providecommand{\Sp}{\operatorname{Sp}}

% ================================================================ TD : mise en page propre au livret d'exercices
\fancyhf{}
\fancyhead[L]{\footnotesize\sffamily\color{Gris}Sélection · Séries numériques}
\fancyhead[R]{\footnotesize\sffamily\scshape\color{Gris}Le Collé}
\fancyfoot[C]{\small\sffamily\color{Ink}\thepage\,/\,\pageref{LastPage}}
\fancypagestyle{plain}{\fancyhf{}\fancyfoot[C]{\small\sffamily\color{Ink}\thepage\,/\,\pageref{LastPage}}\renewcommand{\headrulewidth}{0pt}}
\makeatletter
\newcommand{\partie}[3]{% numéro, titre, description
  \par\needspace{10\baselineskip}\vspace{1.4em}\phantomsection\label{partie:#1}\pdfbookmark[0]{Partie #1}{partie.#1}%
  {\sffamily\small\color{Gris}\textls[200]{PARTIE #1}}\par\nobreak\vspace{0.1em}
  {\sffamily\bfseries\LARGE\color{Ink}#2}\par\nobreak\vspace{0.3em}
  {\color{Ocre}\rule{2.6cm}{1.2pt}}\par\nobreak\vspace{0.2em}
  {\small\color{Gris}#3}\par\nobreak\vspace{0.6em}}
\newcommand{\exercice}[7]{% numéro, difficulté, titre, origine, notions, identifiant, page du corrigé
  \par\needspace{8\baselineskip}\vspace{1.1em}\phantomsection
  \def\@currentlabel{#1}\label{e:#6}\pdfbookmark[1]{Exercice #1}{e.#6}%
  \noindent{\sffamily\bfseries\large\color{Bleu}Exercice~#1}\quad\niveau{#2}\hfill{\footnotesize\sffamily\color{Gris}#4}\par\nobreak
  \vspace{0.1em}\noindent{\sffamily\bfseries\color{Ink}#3}\par\nobreak
  \vspace{0.1em}\noindent{\footnotesize\sffamily\color{Gris}Notions : #5.\hfill\mbox{corrigé p.~#7}\par}\nobreak
  \vspace{0.35em}}
\newcommand{\renvoi}[1]{\par\nopagebreak\vspace{0.2em}\noindent\hfill{\scriptsize\sffamily\color{Gris}#1}\par}
\newcommand{\place}{\par\vspace{1.6cm plus 0.6cm minus 0.6cm}}
\newcommand{\correction}[4]{% numéro, titre, identifiant, renvoi
  \par\needspace{7\baselineskip}\vspace{1.3em}\phantomsection
  \def\@currentlabel{#1}\label{e:#3}\pdfbookmark[1]{Exercice #1}{c.#3}%
  {\color{Filet}\hrule height 0.4pt}\vspace{0.5em}
  \noindent{\sffamily\bfseries\large\color{Bleu}Exercice~#1}\quad{\sffamily\bfseries\color{Ink}#2}\hfill{\scriptsize\sffamily\color{Gris}#4}\par\nobreak
  \vspace{0.4em}}
\newcommand{\sommaireligne}[3]{\noindent\hbox to 2.2em{\hfill\sffamily\color{Bleu}#1.}\enspace\begin{minipage}[t]{0.8\linewidth}\raggedright#2\end{minipage}\hfill\pageref{e:#3}\par}
\makeatother
\setlength{\parskip}{0.5em plus 0.15em minus 0.1em}

\makeatletter
\newcommand{\probleme}[7]{% numéro, étape, titre, difficulté, origine, durée, identifiant
  \clearpage\phantomsection\label{p:#7}\pdfbookmark[1]{Problème #1 — #2}{p.#7}%
  {\sffamily\small\color{Gris}\textls[150]{PROBLÈME #1\ \ ·\ \ \MakeUppercase{#2}}}\par\nobreak\vspace{0.15em}
  {\sffamily\bfseries\LARGE\color{Ink}#3}\par\nobreak\vspace{0.3em}
  {\color{Ocre}\rule{2.6cm}{1.2pt}}\par\nobreak\vspace{0.2em}
  \noindent{\footnotesize\sffamily\color{Gris}Difficulté~\niveau{#4}\quad·\quad durée indicative : #6\quad·\quad #5}\par\nobreak\vspace{0.3em}}
\newcommand{\fichepb}[3]{% notions, objectifs, prérequis
  {\small\setlength{\parskip}{0.15em}%
  \noindent\hangindent=2.4cm\hangafter=1\makebox[2.4cm][l]{\sffamily\textcolor{Bleu}{Notions}}#1\par
  \noindent\hangindent=2.4cm\hangafter=1\makebox[2.4cm][l]{\sffamily\textcolor{Bleu}{Objectifs}}#2\par
  \noindent\hangindent=2.4cm\hangafter=1\makebox[2.4cm][l]{\sffamily\textcolor{Bleu}{Prérequis}}#3\par}
  \vspace{0.4em}{\color{Filet}\hrule height 0.4pt}\vspace{0.6em}}
\newcommand{\titresection}[2]{% titre, texte
  \clearpage\phantomsection\pdfbookmark[0]{#1}{s.#1}%
  {\sffamily\bfseries\fontsize{24}{28}\selectfont\color{Ink}#1}\par\vspace{0.3em}
  {\color{Ocre}\rule{3.2cm}{1.4pt}}\par\vspace{0.5em}{\small\color{Gris}#2}\par\vspace{0.6em}}
\newcommand{\soustitre}[3]{% numéro, titre, label
  \par\needspace{6\baselineskip}\vspace{1em}\phantomsection\label{#3}%
  {\color{Filet}\hrule height 0.4pt}\vspace{0.5em}
  \noindent{\sffamily\bfseries\large\color{Bleu}Problème~#1}\quad{\sffamily\bfseries\color{Ink}#2}\par\nobreak\vspace{0.3em}}
\makeatother

\begin{document}
\thispagestyle{plain}
{\sffamily\small\color{Gris}\textls[150]{LE COLLÉ · MATHÉMATIQUES PT · CHAPITRE 7}}\par\vspace{0.2em}
{\sffamily\bfseries\fontsize{26}{30}\selectfont\color{Ink}Sélection --- Séries}\par\vspace{0.3em}
{\sffamily\Large\color{BleuMoyen}Cinq problèmes à travailler après le cours}\par\vspace{0.4em}
{\color{Ocre}\rule{3.2cm}{1.4pt}}\par\vspace{0.8em}
{\small Ces cinq problèmes sont la sélection principale pour les séances, \emph{après} l'explication du cours.
Ils vont de la mise en application au niveau des devoirs et des concours ; chacun réunit plusieurs résultats du chapitre.
Ils complètent le TD complet et le livre, qu'ils ne remplacent pas.

\textbf{Mode d'emploi.} Les traiter dans l'ordre. Pour chaque problème : lire l'en-tête (notions, objectifs, prérequis),
chercher au brouillon, puis rédiger. En cas de blocage, lire \emph{une seule} indication (page des indications), puis reprendre.
Ne lire la correction qu'une fois le problème terminé ou après un vrai effort : elle explique le \emph{choix de la méthode}
avant les calculs. L'origine de chaque problème est indiquée en tête ; les problèmes « d'après » un concours sont des
adaptations rédigées par Le Collé, pas les énoncés officiels.}
\vspace{1.2em}
{\small\renewcommand{\arraystretch}{1.45}\setlength{\tabcolsep}{4pt}\noindent\begin{tabular}{@{}r>{\raggedright\arraybackslash\sffamily\footnotesize}p{2.9cm}>{\raggedright\arraybackslash}p{6.1cm}ll@{\hspace{6pt}}ccc@{}}
& & & & & {\footnotesize\sffamily énoncé} & {\footnotesize\sffamily indic.} & {\footnotesize\sffamily corrigé}\\\hline
1 & \textcolor{Gris}{Mise en application} & Nature de séries à termes positifs & \niveau{2} & 25\,min & \pageref{p:s-nature} & \pageref{i:s-nature} & \pageref{c:s-nature}\\
2 & \textcolor{Gris}{Méthodes classiques} & Somme d'une série alternée par une intégrale & \niveau{3} & 35\,min & \pageref{p:sn-7} & \pageref{i:sn-7} & \pageref{c:sn-7}\\
3 & \textcolor{Gris}{Approfondissement} & Une série définie par récurrence & \niveau{3} & 40\,min & \pageref{p:sn-11} & \pageref{i:sn-11} & \pageref{c:sn-11}\\
4 & \textcolor{Gris}{Problème de synthèse} & La suite $u_{n+1}=\sin u_n$ & \niveau{4} & 45\,min & \pageref{p:sn-15} & \pageref{i:sn-15} & \pageref{c:sn-15}\\
5 & \textcolor{Gris}{Niveau DS / concours} & Reste d'une série alternée & \niveau{4} & 45\,min & \pageref{p:sn-13} & \pageref{i:sn-13} & \pageref{c:sn-13}\\
\end{tabular}}
\probleme{1}{Mise en application}{Nature de séries à termes positifs}{2}{Livret, exercice 1.5 · TD complet, exercice~7}{25\,min}{s-nature}
\fichepb{équivalents, Riemann, série géométrique, croissances comparées.}{Choisir, pour chaque terme général, un équivalent ou une comparaison qui ramène à une série de référence.}{Développements limités usuels ; Séries de Riemann et géométriques ; Théorème des équivalents (termes positifs).}
Déterminer la nature des séries de terme général :
\[\text{(a)}\ \frac{n^2+1}{n^4+3}\qquad\text{(b)}\ \frac1{\sqrt{n(n+1)}}\qquad\text{(c)}\ 1-\cos\frac1n\qquad\text{(d)}\ \frac1{\ln n+e^n}\qquad\text{(e)}\ n^2e^{-\sqrt n}\qquad\text{(f)}\ \frac{\ln n}n.\]
\par\vspace{0.8em}\noindent\hfill{\scriptsize\sffamily\color{Gris}indications p.~\pageref{i:s-nature} · correction p.~\pageref{c:s-nature}}
\probleme{2}{Méthodes classiques}{Somme d'une série alternée par une intégrale}{3}{Exercice complémentaire · TD complet, exercice~20}{35\,min}{sn-7}
\fichepb{séries alternées, série géométrique, intégration, éléments simples.}{Prouver la convergence par le critère, puis calculer la somme en écrivant chaque terme comme une intégrale.}{Critère des séries alternées ; Somme géométrique finie ; Primitives de $\frac1{1+t}$ et $\frac1{t^2+a^2}$.}
\begin{enumerate}
\item Montrer que $\sum\frac{(-1)^n}{3n+1}$ converge.
\item Montrer que, pour tout $n$, $\displaystyle\sum_{k=0}^n\frac{(-1)^k}{3k+1}=\int_0^1\frac{dt}{1+t^3}-\int_0^1\frac{(-t^3)^{n+1}}{1+t^3}\,dt$.
\item En déduire $\displaystyle\sum_{n=0}^{+\infty}\frac{(-1)^n}{3n+1}=\int_0^1\frac{dt}{1+t^3}$, puis calculer cette intégrale
(décomposer $\frac1{1+t^3}=\frac\alpha{1+t}+\frac{\beta t+\gamma}{t^2-t+1}$).
\end{enumerate}
\par\vspace{0.8em}\noindent\hfill{\scriptsize\sffamily\color{Gris}indications p.~\pageref{i:sn-7} · correction p.~\pageref{c:sn-7}}
\probleme{3}{Approfondissement}{Une série définie par récurrence}{3}{Exercice complémentaire · TD complet, exercice~21}{40\,min}{sn-11}
\fichepb{suite récurrente, série télescopique, équivalents, signe constant.}{Étudier une suite récurrente puis la nature de séries associées, en utilisant le lien entre une suite et la série de ses différences.}{Théorème de la limite monotone ; Séries télescopiques ; Théorème des équivalents.}
Soit $u_0\in\,]-1,1[$ et $u_{n+1}=\frac{1+u_n^2}2$. On pose $w_n=1-u_n$.
\begin{enumerate}
\item Montrer que $u_n\in[0,1[$ pour $n\geq1$, que $(u_n)$ est croissante, et qu'elle tend vers $1$.
\item Montrer que $w_{n+1}=w_n-\frac{w_n^2}2$. En déduire que $\sum w_n^2$ converge.
\item Déterminer la nature de $\sum\ln\frac{w_{n+1}}{w_n}$, puis celle de $\sum w_n$.
\end{enumerate}
\par\vspace{0.8em}\noindent\hfill{\scriptsize\sffamily\color{Gris}indications p.~\pageref{i:sn-11} · correction p.~\pageref{c:sn-11}}
\probleme{4}{Problème de synthèse}{La suite $u_{n+1}=\sin u_n$}{4}{Exercice complémentaire · TD complet, exercice~25}{45\,min}{sn-15}
\fichepb{suite récurrente, lemme de Cesàro, développement limité, équivalent, Riemann.}{Obtenir un équivalent d'une suite récurrente par la méthode classique « $\frac1{u_{n+1}^\alpha}-\frac1{u_n^\alpha}$ a une limite », puis en déduire la nature de séries.}{Développement limité de $\sin$ ; Définition de la limite avec $\varepsilon$ ; Séries de Riemann.}
Soit $u_0\in\,]0,\pi[$ et $u_{n+1}=\sin u_n$.
\begin{enumerate}
\item Montrer que $u_n\in\,]0,\pi[$, que $(u_n)$ décroît et tend vers $0$.
\item \emph{Lemme de Cesàro.} Montrer que si $x_n\to\ell\in\R$, alors $\frac1n\sum_{k=1}^nx_k\to\ell$.
\item Montrer que $\frac1{u_{n+1}^2}-\frac1{u_n^2}\to\frac13$. En déduire que $u_n\sim\sqrt{\frac3n}$.
\item Nature de $\sum u_n$, $\sum u_n^2$, $\sum u_n^3$ ?
\end{enumerate}
\par\vspace{0.8em}\noindent\hfill{\scriptsize\sffamily\color{Gris}indications p.~\pageref{i:sn-15} · correction p.~\pageref{c:sn-15}}
\probleme{5}{Niveau DS / concours}{Reste d'une série alternée}{4}{Exercice complémentaire · TD complet, exercice~23}{45\,min}{sn-13}
\fichepb{reste, séries alternées, développement asymptotique, semi-convergence.}{Obtenir un équivalent du reste d'une série alternée, puis étudier la série des restes : le type de raisonnement des épreuves écrites.}{Critère des séries alternées et majoration du reste ; Développement asymptotique d'un terme de signe variable.}
Soit $R_n=\sum_{k=n+1}^{+\infty}\frac{(-1)^k}k$ le reste de la série $\sum_{k\geq1}\frac{(-1)^k}k$.
\begin{enumerate}
\item Calculer $R_n-R_{n+1}$. Montrer que $R_n+R_{n+1}=\sum_{k=n+1}^{+\infty}\frac{(-1)^k}{k(k+1)}$ et que $|R_n+R_{n+1}|\leq\frac1{(n+1)(n+2)}$.
\item En déduire que $R_n\sim\frac{(-1)^{n+1}}{2n}$, puis que la série $\sum R_n$ converge.
\end{enumerate}
\par\vspace{0.8em}\noindent\hfill{\scriptsize\sffamily\color{Gris}indications p.~\pageref{i:sn-13} · correction p.~\pageref{c:sn-13}}
\titresection{Indications}{Une indication par question délicate. N\'en lire qu\'une à la fois, puis reprendre la recherche.}
\soustitre{1}{Nature de séries à termes positifs}{i:s-nature}
\begin{itemize}[leftmargin=1.2em,itemsep=0.25em]\item (a)(b)(c) Garder les termes dominants ; $1-\cos x\sim\frac{x^2}2$.
\item (d) $\ln n=o(e^n)$ : quel est le terme dominant du dénominateur ?
\item (e) Montrer que $n^2\cdot n^2e^{-\sqrt n}\to0$ en passant à l'exponentielle.
\item (f) Comparer $\frac{\ln n}n$ à $\frac1n$.
\end{itemize}
\soustitre{2}{Somme d'une série alternée par une intégrale}{i:sn-7}
\begin{itemize}[leftmargin=1.2em,itemsep=0.25em]\item 1. Vérifier les trois hypothèses du critère.
\item 2. $\frac{(-1)^k}{3k+1}=\int_0^1(-t^3)^k\,dt$ ; sommer la progression géométrique.
\item 3. Majorer le reste intégral par $\int_0^1t^{3n+3}\,dt$ ; pour l'intégrale, mettre $t^2-t+1$ sous forme canonique.
\end{itemize}
\soustitre{3}{Une série définie par récurrence}{i:sn-11}
\begin{itemize}[leftmargin=1.2em,itemsep=0.25em]\item 1. $u_{n+1}-u_n=\frac{(1-u_n)^2}2$.
\item 2. Écrire $w_n^2$ comme $2(w_n-w_{n+1})$.
\item 3. $\ln\frac{w_{n+1}}{w_n}=\ln w_{n+1}-\ln w_n$, et $\ln\left(1-\frac{w_n}2\right)\sim-\frac{w_n}2$.
\end{itemize}
\soustitre{4}{La suite $u_{n+1}=\sin u_n$}{i:sn-15}
\begin{itemize}[leftmargin=1.2em,itemsep=0.25em]\item 1. $\sin x\leq x$ sur $\R_+$ ; limite : point fixe de $\sin$.
\item 2. Couper la somme en $k\leq N$ et $k>N$.
\item 3. Développer $\sin^2x$ à l'ordre 4 ; appliquer Cesàro aux différences.
\item 4. Élever l'équivalent aux puissances $1$, $2$, $3$.
\end{itemize}
\soustitre{5}{Reste d'une série alternée}{i:sn-13}
\begin{itemize}[leftmargin=1.2em,itemsep=0.25em]\item 1. Écrire $R_{n+1}$ avec l'indice décalé, puis regrouper terme à terme.
\item 2. $2R_n=(R_n-R_{n+1})+(R_n+R_{n+1})$.
\item 3. Un équivalent ne suffit pas pour un signe variable : garder le $O\left(\frac1{n^2}\right)$.
\end{itemize}
\titresection{Corrections}{Chaque correction commence par le choix de la méthode ; les encadrés résument ce qui resservira. Les renvois « méthode~1.4 du livret » désignent le livret du chapitre.}
\soustitre{1}{Nature de séries à termes positifs}{c:s-nature}
Tous les termes sont positifs.
\q{(a)} $\frac{n^2+1}{n^4+3}\sim\frac1{n^2}$ : convergente (Riemann, $2>1$).
\q{(b)} $\frac1{\sqrt{n(n+1)}}\sim\frac1n$ : divergente.
\q{(c)} $1-\cos\frac1n\sim\frac1{2n^2}$ : convergente.
\q{(d)} $\ln n=o(e^n)$, donc $\frac1{\ln n+e^n}\sim e^{-n}=\left(\frac1e\right)^n$ : convergente (géométrique de raison $\frac1e<1$).
\q{(e)} $n^2\cdot n^2e^{-\sqrt n}=e^{4\ln n-\sqrt n}\to0$ car $\ln n=o(\sqrt n)$ ; donc $n^2e^{-\sqrt n}=o\left(\frac1{n^2}\right)$ : convergente.
\q{(f)} $\frac{\ln n}n\geq\frac1n$ pour $n\geq3$ : divergente par comparaison.
\par\hfill\textcolor{BleuClair}{$\square$}\par
\soustitre{2}{Somme d'une série alternée par une intégrale}{c:sn-7}
\q{1.} $\frac1{3n+1}$ est positive, décroissante, de limite nulle : critère des séries alternées.
\q{2.} $\sum_{k=0}^n(-t^3)^k=\frac{1-(-t^3)^{n+1}}{1+t^3}$ ; on intègre sur $[0,1]$, avec $\int_0^1(-1)^kt^{3k}\,dt=\frac{(-1)^k}{3k+1}$.
\q{3.} $\left|\int_0^1\frac{(-t^3)^{n+1}}{1+t^3}\,dt\right|\leq\int_0^1t^{3n+3}\,dt=\frac1{3n+4}\to0$. Ensuite $\alpha=\frac13$, $\beta=-\frac13$,
$\gamma=\frac23$ : $\frac{-t+2}{t^2-t+1}=-\frac12\cdot\frac{2t-1}{t^2-t+1}+\frac32\cdot\frac1{(t-\frac12)^2+\frac34}$. D'où
\[\int_0^1\frac{dt}{1+t^3}=\frac13\ln2-\frac16\big[\ln(t^2-t+1)\big]_0^1+\frac12\cdot\frac2{\sqrt3}\left[\arctan\frac{2t-1}{\sqrt3}\right]_0^1=\frac{\ln2}3+\frac\pi{3\sqrt3}.\]
(Numériquement $0{,}8356\dots$, vérifié.)
\par\hfill\textcolor{BleuClair}{$\square$}\par
\soustitre{3}{Une série définie par récurrence}{c:sn-11}
\q{1.} $u_1=\frac{1+u_0^2}2\in\left[\frac12,1\right[$, et si $u_n\in[0,1[$, alors $u_{n+1}\in\left[\frac12,1\right[$. Ensuite
$u_{n+1}-u_n=\frac{(1-u_n)^2}2\geq0$. Croissante et majorée par $1$, $(u_n)$ converge vers $\ell$ avec $\ell=\frac{1+\ell^2}2$, soit $\ell=1$.
\q{2.} $w_{n+1}=1-\frac{1+(1-w_n)^2}2=w_n-\frac{w_n^2}2$. Donc $w_n^2=2(w_n-w_{n+1})$ : série télescopique, convergente car $(w_n)$ converge
(vers $0$), de somme $2w_0$.
\q{3.} $\ln\frac{w_{n+1}}{w_n}=\ln w_{n+1}-\ln w_n$ et $\ln w_n\to-\infty$ : la série télescopique diverge. Or
$\ln\frac{w_{n+1}}{w_n}=\ln\left(1-\frac{w_n}2\right)\sim-\frac{w_n}2$, de signe constant (négatif) : $\sum w_n$ a la même nature, elle
\textbf{diverge}.
\retenir{une suite $(a_n)$ et la série $\sum(a_{n+1}-a_n)$ ont même nature : c'est un outil pour étudier des suites définies par récurrence.}
\par\hfill\textcolor{BleuClair}{$\square$}\par
\soustitre{4}{La suite $u_{n+1}=\sin u_n$}{c:sn-15}
\q{1.} Si $u_n\in\,]0,\pi[$, alors $u_{n+1}=\sin u_n\in\,]0,1]\subset\,]0,\pi[$. Comme $\sin x\leq x$ sur $\R_+$, $u_{n+1}\leq u_n$ : la suite décroît,
minorée par $0$, elle converge vers $\ell\in[0,1]$ avec $\sin\ell=\ell$, soit $\ell=0$.
\q{2.} Soit $\varepsilon>0$ et $N$ tel que $|x_k-\ell|\leq\varepsilon$ pour $k>N$. Pour $n>N$ :
$\left|\frac1n\sum_{k=1}^n(x_k-\ell)\right|\leq\frac1n\sum_{k=1}^N|x_k-\ell|+\frac{n-N}n\varepsilon$. Le premier terme tend vers $0$ ($N$ fixé), donc
il est $\leq\varepsilon$ pour $n$ grand : la moyenne est à moins de $2\varepsilon$ de $\ell$.
\q{3.} Avec $x=u_n\to0$ : $\frac1{\sin^2x}-\frac1{x^2}=\frac{x^2-\sin^2x}{x^2\sin^2x}$, et $\sin^2x=x^2-\frac{x^4}3+o(x^4)$, donc le quotient vaut
$\frac{x^4/3+o(x^4)}{x^4+o(x^4)}\to\frac13$. Par Cesàro, $\frac1n\sum_{k=0}^{n-1}\left(\frac1{u_{k+1}^2}-\frac1{u_k^2}\right)=\frac1n\left(\frac1{u_n^2}-\frac1{u_0^2}\right)\to\frac13$,
donc $\frac1{u_n^2}\sim\frac n3$ et $u_n\sim\sqrt{\frac3n}$.
\q{4.} Termes positifs : $u_n\sim\frac{\sqrt3}{n^{1/2}}$ (divergente), $u_n^2\sim\frac3n$ (divergente), $u_n^3\sim\frac{3\sqrt3}{n^{3/2}}$ (convergente).
\par\hfill\textcolor{BleuClair}{$\square$}\par
\soustitre{5}{Reste d'une série alternée}{c:sn-13}
\q{1.} $R_n-R_{n+1}=\frac{(-1)^{n+1}}{n+1}$. Ensuite $R_{n+1}=\sum_{k=n+1}^{+\infty}\frac{(-1)^{k+1}}{k+1}$ (changement d'indice), donc
$R_n+R_{n+1}=\sum_{k\geq n+1}(-1)^k\left(\frac1k-\frac1{k+1}\right)=\sum_{k\geq n+1}\frac{(-1)^k}{k(k+1)}$. C'est le reste d'une série alternée dont le
terme $\frac1{k(k+1)}$ décroît vers $0$ : il est majoré en valeur absolue par son premier terme $\frac1{(n+1)(n+2)}$.
\q{2.} $2R_n=(R_n-R_{n+1})+(R_n+R_{n+1})=\frac{(-1)^{n+1}}{n+1}+O\left(\frac1{n^2}\right)$, donc $R_n=\frac{(-1)^{n+1}}{2(n+1)}+O\left(\frac1{n^2}\right)$ et
$R_n\sim\frac{(-1)^{n+1}}{2n}$. La série $\sum R_n$ est la somme d'une série alternée convergente et d'une série absolument convergente :
elle converge.
\piege{conclure « $\sum R_n$ converge car $R_n\sim\frac{(-1)^{n+1}}{2n}$ » : l'équivalent ne suffit pas pour une série de signe variable. C'est le
développement $\frac{(-1)^{n+1}}{2(n+1)}+O(\frac1{n^2})$ qui permet de conclure.}
\par\hfill\textcolor{BleuClair}{$\square$}\par
\end{document}
